\chapter{Fundamentação Teórica}
\label{chap:fundamentacao}

A presente análise teórica explora os conceitos que fundamentam a coexistência entre JavaScript e TypeScript no desenvolvimento de aplicações web. O percurso argumentativo parte da tipagem dinâmica e da assincronicidade do JavaScript, apresenta o TypeScript como uma forma de tipagem gradual e discute os limites dessa proteção. Em seguida, examina contratos de bibliotecas, configuração do compilador, ferramentas de migração e a mudança do perfil de falhas em ecossistemas modernos, posicionando a experiência do desenvolvedor como um fator relevante, mas não isolado, na escolha tecnológica.

O desenvolvimento front-end foi profundamente moldado pelo papel central do \textit{JavaScript}, linguagem criada em 1995 por Brendan Eich. Embora concebido inicialmente como um recurso de script para adicionar interatividade a páginas estáticas, sua evolução, marcada pela padronização como \textit{ECMAScript}, o consolidou como a principal tecnologia para aplicações web, garantindo consistência entre diferentes ambientes de execução e tornando-se o padrão universal na camada de interação com o usuário \cite{mororo2024}. Uma de suas características mais marcantes é a tipagem dinâmica, na qual os tipos de dados estão associados aos valores, e não às variáveis. Essa flexibilidade permite uma escrita de código ágil, tornando a linguagem atrativa para prototipagem e projetos menores, onde uma mesma variável pode mudar de tipo dinamicamente sem que o desenvolvedor precise de declarações prévias.

\begin{verbatim}
let resultado = 100; // 'resultado' é inferido como number
resultado = 'cem';   // Válido em JS. Agora 'resultado' é uma string.
\end{verbatim}

Contudo, essa mesma liberdade introduz uma classe de desafios em sistemas complexos. Como aponta \cite{makimov2022}, erros de tipo, como a passagem de dados em formato inesperado para uma função, podem só se manifestar em tempo de execução. Tal característica eleva o custo de manutenção, pois a depuração de erros em grandes bases de código torna-se mais árdua. A essa complexidade soma-se o modelo de concorrência do JavaScript, baseado em um \textit{event loop}. Embora eficaz para impedir que operações demoradas bloqueiem a interface, a gestão do fluxo assíncrono permanece uma fonte de condições de corrida, falhas de tratamento e comportamentos difíceis de reproduzir. A introdução de \textit{Promises} e, posteriormente, da sintaxe \textit{async/await}, melhorou a legibilidade, mas não substituiu a necessidade de testes e validações em execução \cite{mororo2024,tang2026ecosystem}.

Foi nesse cenário de um ecossistema poderoso, porém propenso a desafios de escalabilidade, que a Microsoft lançou o \textit{TypeScript} em 2012. Definido como um superconjunto de JavaScript, sua principal inovação foi a introdução de um sistema de tipagem estática opcional. Essa abordagem significa que todo código JavaScript é, por definição, um código TypeScript válido, o que viabiliza uma migração gradual e estratégica. A proposta central do TypeScript é deslocar a verificação de tipos do tempo de execução para o tempo de compilação, proporcionando maior segurança e melhorando a legibilidade, visto que o código se torna, em parte, autodescritivo \cite{george2020, makimov2022}. Um erro de atribuição de tipo, por exemplo, é capturado antes mesmo que o programa seja executado.

\begin{verbatim}
let meuNome: string = "Gustavo";
// O compilador do TypeScript apontará um erro nesta linha:
// Erro: O tipo 'number' não pode ser atribuído ao tipo 'string'.
meuNome = 30; 
\end{verbatim}

A promessa de maior robustez do TypeScript deve, portanto, ser interpretada de forma condicional. Um estudo de \cite{richards2017} estimou que a tipagem estática poderia prevenir cerca de 15\% dos bugs públicos em projetos JavaScript. Por outro lado, a análise de \cite{bogner2022} indicou que aplicações TypeScript podem apresentar melhor qualidade de código sem necessariamente possuir menos bugs. Esse resultado é compatível com a ideia de que a tipagem reduz determinadas falhas locais, mas não substitui a análise da lógica, dos requisitos e das integrações.

Essa distinção é reforçada por estudos sobre o próprio TypeScript. Wang, Fang e Wang \cite{wang2025bugs} identificaram uma variedade de defeitos ligados a nomes, bindings, escopos, tipos, semântica, recursos da linguagem e comportamento do compilador. Em uma escala mais ampla, Tang, Alimadadi e Sumner \cite{tang2026ecosystem} analisaram 633 relatos de bugs em 16 projetos e observaram a predominância de problemas de configuração e ferramentas, uso incorreto de APIs e tratamento assíncrono. Os resultados não autorizam concluir que o TypeScript seja menos confiável; indicam que a adoção desloca parte da atenção necessária para as camadas de integração e orquestração.

Para além da tipagem, o TypeScript oferece interfaces, tipos genéricos e mecanismos de composição que podem funcionar como documentação e apoio à refatoração. Entretanto, uma declaração de tipo só é útil quando representa o comportamento real do componente. Feldthaus e Møller \cite{feldthaus2014interfaces} verificaram a correção de interfaces TypeScript para bibliotecas JavaScript e demonstraram a importância de testar o contrato contra a implementação. De modo semelhante, Hoeflich, Findler e Serrano \cite{hoeflich2022scotty} desenvolveram o Scotty para encontrar incompatibilidades no repositório DefinitelyTyped, mostrando que declarações incorretas podem produzir uma sensação indevida de segurança. A fronteira entre o tipo declarado e o comportamento efetivo da biblioteca é, portanto, um elemento central da confiabilidade.

Apesar dos benefícios, a adoção do TypeScript impõe uma curva de aprendizado e custos de configuração. Desenvolvedores precisam compreender novos conceitos, integrar compilador, bundler, testes e linters e decidir quanto rigor aplicar ao projeto. Em estudo com mais de 4.000 repositórios, Zandonotto et al. \cite{zandonotto2026compiler} encontraram associação entre opções como \texttt{noImplicitAny} e \texttt{noUnusedParameters} e menores índices de problemas estáticos, mas também observaram correlações contraintuitivas para algumas opções. Como o estudo é correlacional, seus resultados não constituem garantia causal; ainda assim, evidenciam que a simples presença de arquivos \texttt{.ts} não define o nível de proteção obtido.

A migração pode ser apoiada por ferramentas de inferência e geração automática. O dataset ManyTypes4TypeScript reúne projetos e anotações de tipos em escala para pesquisas de inferência \cite{jesse2022manytypes}, enquanto Jesse, Devanbu e Sawant \cite{jesse2023types} investigaram a previsão de tipos definidos por usuários. Cristiani e Thiemann \cite{cristiani2021declarations} estudaram a geração de arquivos de declaração a partir de JavaScript. Essas técnicas podem reduzir o custo inicial, mas não eliminam a revisão humana: uma inferência ou declaração incorreta pode propagar um contrato falso pelo sistema.

Como aponta \cite{shrestha2022}, o conhecimento prévio de JavaScript também pode interferir na adaptação ao novo paradigma. Ainda assim, a integração com frameworks e ferramentas modernas e os ganhos percebidos em projetos de grande porte ajudam a explicar a relevância do TypeScript \cite{stackoverflow2023}. A literatura, portanto, não sustenta uma oposição simples entre uma linguagem melhor e outra pior. O JavaScript permanece vantajoso para aprendizagem, prototipagem e situações em que a agilidade inicial é prioritária; o TypeScript tende a oferecer maior apoio à manutenção e à coordenação em bases extensas, desde que acompanhado por contratos confiáveis, configuração coerente, testes e validações de runtime. Esse equilíbrio orienta a análise empírica deste trabalho.
